% SPDX-License-Identifier: CC-BY-SA-4.0
% Author: Matthieu Perrin
% Part: <Nom de la partie>
% Section: <Nom de la section>
% Sub-section: <Nom de la sous-section>  % (facultatif, laisser vide si non utilisé)
% Frame: <Titre de la slide>

\begingroup

\begin{frame}{Vers une axiomatisation nécessairement incomplète}
 
  \onBlock[right=.75\textwidth, top=-2mm]{Théorème d'incomplétude de Gödel (1931)}{
    Il n'existe pas de système de preuve :
    \begin{itemize}
    \item formel (\textsc{Vérification des preuves} est décidable)
    \item suffisamment expressif pour l'arithmétique entière
    \item à la fois complet et cohérent
    \end{itemize}
  }
 
  \onBlock<2->[left=.73\textwidth, bottom=-3mm]{Vers une théorie des ensembles cohérente}{
    \begin{itemize}
    \item Schéma de compréhension restreint
      $\structure{\{x \alert{\in E} \mid P(x) \}}$
    \item La base de l'axiomatisation moderne (ZFC)
      \begin{itemize}
      \item Complétée par Fraenkel, Skolem, Von Neumann...
      \end{itemize}
    \item Cohérence toujours indémontrée
    \end{itemize}
    \begin{minipage}{.92\textwidth}
      \myquote{David Hilbert}{Personne ne doit pouvoir nous chasser\\ du paradis que Cantor nous a créé}
    \end{minipage}
  }
  
  \onImage[x=-.4\textwidth,top]{%
    height=25mm,
    title={Kurt Gödel},
    licenselogo={\ccPublicDomain},
    license*={Domaine public (Vienne, 1925. \href{https://commons.wikimedia.org/wiki/File:Young_Kurt_G\%C3\%B6del_as_a_student_in_1925.jpg}{Wikimedia})},
    img={Godel.jpg}
  }
 
  \onImage<2->[x=.4\textwidth,bottom=5mm]{%
    height=30mm,
    title={Ernst Zermelo},
    licenselogo={\ccPublicDomain},
    license*={Domaine public (Italie, vers 1900. \href{https://commons.wikimedia.org/wiki/File:Ernst_Zermelo_1900s.jpg}{Wikimedia})},
    img={Zermelo.jpg}
  }
  
\end{frame}

\endgroup
